\documentclass[11pt]{article}
\usepackage[utf8]{inputenc}
\usepackage{amsmath,amssymb,booktabs,geometry,verbatim}
\usepackage{CJKutf8}
\geometry{margin=1in}
\newcommand{\lbr}{\char`\{}
\newcommand{\rbr}{\char`\}}

\title{CSP-J 2023 (Junior)}
\date{2023-10-21, 8:30--12:00 (3.5 hours)}

\begin{document}
\maketitle

\noindent
Second round of the 2023 CCF Certified Software Professional (CSP-J/S 2023),
junior level (CSP-J).
Four problems; each problem is worth 100 points.
In the original contest every problem reads its input from a file
(\texttt{apple.in}, \texttt{road.in}, \texttt{uqe.in}, \texttt{bus.in})
and writes its output to a file
(\texttt{apple.out}, \texttt{road.out}, \texttt{uqe.out}, \texttt{bus.out});
the sentences stating this are kept in the translation below.
Time limit and memory limit are those shown on QOJ.

\section*{Problem 1: Small Apples (\begin{CJK}{UTF8}{gbsn}小苹果\end{CJK})}

\noindent
Time limit: 1 s.
Memory limit: 512 MB.
Total points: 100.
Code name: \texttt{apple}.

\subsection*{Problem description}

There are $n$ apples on Little Y's table, arranged in a row from left to
right and numbered from $1$ to $n$.
Little Bao is a good friend of Little Y, and every day she takes away some of
the apples.
Every day, when taking apples, Little Bao starts from the first apple on the
left and takes away one apple after every two apples that she skips
(that is, she takes the first apple, then skips two, takes one, and so on).
Afterwards, Little Bao arranges the remaining apples into a row again in
their original order.
Little Bao wants to know how many days it takes to take away all the apples,
and on which day the apple numbered $n$ is taken away.

\subsection*{Input}

Read the input from the file \texttt{apple.in}.
The first line of the input contains one positive integer $n$,
the total number of apples.

\subsection*{Output}

Write the output to the file \texttt{apple.out}.
Output one line containing two positive integers separated by a single
space:
the number of days Little Bao needs to take away all the apples,
and the day on which the apple numbered $n$ is taken away.

\subsection*{Sample 1 input}

\begin{verbatim}
8
\end{verbatim}

\subsection*{Sample 1 output}

\begin{verbatim}
5 5
\end{verbatim}

\subsection*{Explanation of sample 1}

There are $8$ apples on the table in total.
On the first day, Little Bao takes away the apples numbered $1$, $4$, $7$.
On the second day, Little Bao takes away the apples numbered $2$, $6$.
On the third day, Little Bao takes away the apple numbered $3$.
On the fourth day, Little Bao takes away the apple numbered $5$.
On the fifth day, Little Bao takes away the apple numbered $8$.

\subsection*{Sample 2}

See the files \texttt{apple/apple2.in} and \texttt{apple/apple2.ans}
in the contestant's directory.
(This sample is given only as downloadable files; it is not printed in the
statement.)

\subsection*{Constraints}

For all test data: $1 \le n \le 10^9$.

\begin{center}
\begin{tabular}{ccc}
\toprule
Test cases & $n \le$ & Special property \\
\midrule
$1 \sim 2$ & $10$ & None \\
$3 \sim 5$ & $10^3$ & None \\
$6 \sim 7$ & $10^6$ & Yes \\
$8 \sim 9$ & $10^6$ & None \\
$10$ & $10^9$ & None \\
\bottomrule
\end{tabular}
\end{center}

Special property: Little Bao takes away the apple numbered $n$ on the first
day.

\section*{Problem 2: Highway (\begin{CJK}{UTF8}{gbsn}公路\end{CJK})}

\noindent
Time limit: 1 s.
Memory limit: 512 MB.
Total points: 100.
Code name: \texttt{road}.

\subsection*{Problem description}

Little Bao is going on a self-driving trip along a highway.
There are $n$ stations on the highway, numbered from $1$ to $n$.
The distance between station $i$ and station $i+1$ is $v_i$ kilometres.
Fuel can be bought at every station on the highway.
At the station numbered $i$ one litre of fuel costs $a_i$ yuan,
and every station sells only whole litres of fuel.
Little Bao wants to drive from station $1$ to station $n$.
At the beginning Little Bao is at station $1$ and the fuel tank of the car is
empty.
It is known that the fuel tank is large enough to hold any amount of fuel,
and that every litre of fuel lets the car travel $d$ kilometres.
What is the minimum amount of money Little Bao has to spend on fuel to drive
from station $1$ to station $n$?

\subsection*{Input}

Read the input from the file \texttt{road.in}.
The first line of the input contains two positive integers $n$ and $d$,
the number of stations on the highway and the distance the car can travel per
litre of fuel.
The second line of the input contains $n-1$ positive integers
$v_1, v_2, \dots, v_{n-1}$, the distances between consecutive stations.
The third line of the input contains $n$ positive integers
$a_1, a_2, \dots, a_n$, the fuel prices at the stations.
(The original statement says ``the second line'' for both of the last two
lines; the prices are on the third line.)

\subsection*{Output}

Write the output to the file \texttt{road.out}.
Output one line containing a single positive integer:
the minimum amount of money Little Bao has to spend on fuel to drive from
station $1$ to station $n$.

\subsection*{Sample 1 input}

\begin{verbatim}
5 4
10 10 10 10
9 8 9 6 5
\end{verbatim}

\subsection*{Sample 1 output}

\begin{verbatim}
79
\end{verbatim}

\subsection*{Explanation of sample 1}

In an optimal plan, Little Bao buys $3$ litres of fuel at station $1$,
$5$ litres of fuel at station $2$, and $2$ litres of fuel at station $4$.

\subsection*{Sample 2}

See the files \texttt{road/road2.in} and \texttt{road/road2.ans}
in the contestant's directory.
(This sample is given only as downloadable files; it is not printed in the
statement.)

\subsection*{Constraints}

For all test data it is guaranteed that
$1 \le n \le 10^5$, $1 \le d \le 10^5$, $1 \le v_i \le 10^5$,
$1 \le a_i \le 10^5$.

\begin{center}
\begin{tabular}{ccc}
\toprule
Test cases & $n \le$ & Special property \\
\midrule
$1 \sim 5$ & $8$ & None \\
$6 \sim 10$ & $10^3$ & None \\
$11 \sim 13$ & $10^5$ & A \\
$14 \sim 16$ & $10^5$ & B \\
$17 \sim 20$ & $10^5$ & None \\
\bottomrule
\end{tabular}
\end{center}

Special property A: the fuel price at station $1$ is the lowest.
Special property B: for all $1 \le i < n$, $v_i$ is a multiple of $d$.

\section*{Problem 3: Quadratic Equation
  (\begin{CJK}{UTF8}{gbsn}一元二次方程\end{CJK})}

\noindent
Time limit: 1 s.
Memory limit: 512 MB.
Total points: 100.
Code name: \texttt{uqe}.

\subsection*{Background}

As is well known, the real solutions of a quadratic equation
$ax^2 + bx + c = 0$ ($a \ne 0$) can be found as follows.
\begin{itemize}
\item Compute $\Delta = b^2 - 4ac$; then:
  \begin{enumerate}
  \item if $\Delta < 0$, the quadratic equation has no real solution;
  \item otherwise $\Delta \ge 0$, and the quadratic equation has the two
    real solutions
    \[
      x_{1,2} = \frac{-b \pm \sqrt{\Delta}}{2a};
    \]
    \begin{itemize}
    \item here $\sqrt{\Delta}$ denotes the arithmetic square root of
      $\Delta$, that is, the unique non-negative real number $s$ with
      $s^2 = \Delta$;
    \item in particular, when $\Delta = 0$ the two real solutions are
      equal, and when $\Delta > 0$ the two real solutions are distinct.
    \end{itemize}
  \end{enumerate}
\end{itemize}
For example:
\begin{itemize}
\item $x^2 + x + 1 = 0$ has no real solution, because
  $\Delta = 1^2 - 4 \times 1 \times 1 = -3 < 0$;
\item $x^2 - 2x + 1 = 0$ has two equal real solutions $x_{1,2} = 1$;
\item $x^2 - 3x + 2 = 0$ has two distinct real solutions $x_1 = 1$,
  $x_2 = 2$.
\end{itemize}
In the statement, the greatest common divisor of $a$ and $b$ is denoted by
$\gcd(a, b)$.
For example, the greatest common divisor of $12$ and $18$ is $6$,
that is, $\gcd(12, 18) = 6$.

\subsection*{Problem description}

You are given the coefficients $a, b, c$ of a quadratic equation, where
$a, b, c$ are \emph{all integers and $a \ne 0$}.
You need to decide whether the quadratic equation $ax^2 + bx + c = 0$ has
real solutions, and print the answer in the required format.

\emph{In this problem, when a rational number $v$ is printed, the following
rules must be followed}:
\begin{itemize}
\item by the definition of a rational number, there exist \emph{unique}
  integers $p$ and $q$ with $q > 0$, $\gcd(p, q) = 1$ and $v = \frac{p}{q}$;
\item \emph{if $q = 1$, print} \texttt{\lbr{}p\rbr{}}; \emph{otherwise print}
  \texttt{\lbr{}p\rbr{}/\lbr{}q\rbr{}},
  where \texttt{\lbr{}n\rbr{}} stands for the value of the integer $n$;
\item for example:
  \begin{itemize}
  \item when $v = -0.5$, the values of $p$ and $q$ are $-1$ and $2$,
    so \texttt{-1/2} should be printed;
  \item when $v = 0$, the values of $p$ and $q$ are $0$ and $1$,
    so \texttt{0} should be printed.
  \end{itemize}
\end{itemize}

\emph{For solving the equation, two cases are distinguished}:
\begin{enumerate}
\item If $\Delta = b^2 - 4ac < 0$, the equation has no real solution,
  and you should print \texttt{NO}.
\item Otherwise $\Delta \ge 0$, and the equation has two solutions
  (possibly equal); let the \emph{larger} one be $x$. Then:
  \begin{enumerate}
  \item[(1)] if $x$ is a rational number, print $x$ in the format for
    rational numbers;
  \item[(2)] otherwise, by the formula above, $x$ can be written
    \emph{uniquely} in the form $x = q_1 + q_2 \sqrt{r}$, where
    \begin{itemize}
    \item $q_1, q_2$ are rational numbers and $q_2 > 0$;
    \item $r$ is a positive integer with $r > 1$, and there is no positive
      integer $d > 1$ with $d^2 \mid r$ (that is, $r$ must not be a
      multiple of $d^2$);
    \end{itemize}
    in this case:
    \begin{enumerate}
    \item if $q_1 \ne 0$, print $q_1$ in the format for rational numbers,
      followed by a plus sign \texttt{+};
    \item otherwise skip this step;
    \end{enumerate}
    and then:
    \begin{enumerate}
    \item if $q_2 = 1$, print \texttt{sqrt(\lbr{}r\rbr{})};
    \item otherwise, if $q_2$ is an integer, print
      \texttt{\lbr{}q2\rbr{}*sqrt(\lbr{}r\rbr{})};
    \item otherwise, if $q_3 = \frac{1}{q_2}$ is an integer, print
      \texttt{sqrt(\lbr{}r\rbr{})/\lbr{}q3\rbr{}};
    \item otherwise, it can be shown that there exist unique integers
      $c, d$ with $c, d > 1$, $\gcd(c, d) = 1$ and $q_2 = \frac{c}{d}$;
      in this case print
      \texttt{\lbr{}c\rbr{}*sqrt(\lbr{}r\rbr{})/\lbr{}d\rbr{}}.
    \end{enumerate}
    In the notation above, \texttt{\lbr{}n\rbr{}} stands for the value of the
    integer $n$; see the samples for details.
  \end{enumerate}
\end{enumerate}
If the equation has real solutions, print the larger of the two real
solutions in the required format.
Otherwise, if the equation has no real solution, print \texttt{NO}.

\subsection*{Input}

Read the input from the file \texttt{uqe.in}.
The first line of the input contains two positive integers $T$ and $M$,
the number of equations and an upper bound on the absolute values of the
coefficients.
Each of the next $T$ lines contains three integers $a, b, c$.

\subsection*{Output}

Write the output to the file \texttt{uqe.out}.
Output $T$ lines, each containing one string: the answer to the
corresponding query, in the format described in the statement.
\emph{The string printed on each line must not contain any space.}

\subsection*{Sample 1 input}

\begin{verbatim}
9 1000
1 -1 0
-1 -1 -1
1 -2 1
1 5 4
4 4 1
1 0 -432
1 -3 1
2 -4 1
1 7 1
\end{verbatim}

\subsection*{Sample 1 output}

\begin{verbatim}
1
NO
1
-1
-1/2
12*sqrt(3)
3/2+sqrt(5)/2
1+sqrt(2)/2
-7/2+3*sqrt(5)/2
\end{verbatim}

\subsection*{Sample 2}

See the files \texttt{uqe/uqe2.in} and \texttt{uqe/uqe2.ans}
in the contestant's directory.
(This sample is given only as downloadable files; it is not printed in the
statement.)

\subsection*{Constraints}

For all test data: $1 \le T \le 5000$, $1 \le M \le 10^3$,
$|a|, |b|, |c| \le M$, $a \ne 0$.

\begin{center}
\begin{tabular}{ccccc}
\toprule
Test case & $M \le$ & Special property A & Special property B &
Special property C \\
\midrule
$1$ & $1$ & Yes & Yes & Yes \\
$2$ & $20$ & No & No & No \\
$3$ & $10^3$ & Yes & No & Yes \\
$4$ & $10^3$ & Yes & No & No \\
$5$ & $10^3$ & No & Yes & Yes \\
$6$ & $10^3$ & No & Yes & No \\
$7, 8$ & $10^3$ & No & No & Yes \\
$9, 10$ & $10^3$ & No & No & No \\
\bottomrule
\end{tabular}
\end{center}

Here:
\begin{itemize}
\item Special property A: it is guaranteed that $b = 0$;
\item Special property B: it is guaranteed that $c = 0$;
\item Special property C: if the equation has solutions, then both solutions
  are integers.
\end{itemize}

\section*{Problem 4: Tour Bus (\begin{CJK}{UTF8}{gbsn}旅游巴士\end{CJK})}

\noindent
Time limit: 1 s.
Memory limit: 512 MB.
Total points: 100.
Code name: \texttt{bus}.

\subsection*{Problem description}

Little Z plans to take a tour bus during the National Day holiday to visit a
scenic area he has long wanted to see.
The map of the scenic area has $n$ places, connected by $m$ roads.
Place number $1$ is the entrance of the scenic area and place number $n$ is
its exit.
We call the time at which the scenic area opens for the day time $0$;
starting from time $0$, every $k$ units of time a tour bus arrives at the
entrance of the scenic area, and at the same moment a tour bus leaves the
scenic area from the exit.
All roads are \emph{one-way}.
For every road, a visitor needs \emph{exactly} $1$ unit of time to walk
through it.
Little Z wants to arrive at the entrance by tour bus, walk to the exit along
any path of his choice, and then leave by tour bus;
this means that both the time he arrives at and the time he leaves the scenic
area must be \emph{non-negative integer multiples} of $k$.
Because of the holiday crowds, \emph{before leaving the scenic area by tour
bus, Little Z wants to keep moving along the roads of the scenic area all the
time, and does not want to stay at any place (including the entrance and the
exit) or on any road.}
Before setting off, Little Z suddenly learns that the scenic area has adopted
a crowd-control measure: every road has been given an ``opening time'' $a_i$,
and visitors may pass through the road only at a time \emph{not earlier than}
$a_i$.
Please help Little Z design a tour plan so that the time at which he leaves
the scenic area by tour bus is as early as possible.

\subsection*{Input}

Read the input from the file \texttt{bus.in}.
The first line of the input contains $3$ positive integers $n, m, k$:
the number of places in the scenic area, the number of roads, and the
interval between tour buses.
Each of the next $m$ lines of the input contains $3$ non-negative integers
$u_i, v_i, a_i$, meaning that the $i$-th road goes from place $u_i$ to place
$v_i$ and has ``opening time'' $a_i$.

\subsection*{Output}

Write the output to the file \texttt{bus.out}.
Output one line containing a single integer: the earliest time at which
Little Z can leave the scenic area by tour bus.
If no tour plan satisfying the requirements exists, output \texttt{-1}.

\subsection*{Sample 1 input}

\begin{verbatim}
5 5 3
1 2 0
2 5 1
1 3 0
3 4 3
4 5 1
\end{verbatim}

\subsection*{Sample 1 output}

\begin{verbatim}
6
\end{verbatim}

\subsection*{Explanation of sample 1}

(The original statement contains a figure, ``Figure 1: illustration of
sample 1'', drawing the five places and the five roads with their opening
times; it is not reproduced here.)
Little Z can arrive at the entrance of the scenic area at time $3$,
walk to the exit in the order $1 \to 3 \to 4 \to 5$,
and leave at time $6$.

\subsection*{Sample 2}

See the files \texttt{bus/bus2.in} and \texttt{bus/bus2.ans}
in the contestant's directory.
(This sample is given only as downloadable files; it is not printed in the
statement.)

\subsection*{Constraints}

For all test data: $2 \le n \le 10^4$, $1 \le m \le 2 \times 10^4$,
$1 \le k \le 100$, $1 \le u_i, v_i \le n$, $0 \le a_i \le 10^6$.

\begin{center}
\begin{tabular}{ccccc}
\toprule
Test case & $n \le$ & $m \le$ & $k \le$ & Special property \\
\midrule
$1 \sim 2$ & $10$ & $15$ & $100$ & $a_i = 0$ \\
$3 \sim 5$ & $10$ & $15$ & $100$ & None \\
$6 \sim 7$ & $10^4$ & $2 \times 10^4$ & $1$ & $a_i = 0$ \\
$8 \sim 10$ & $10^4$ & $2 \times 10^4$ & $1$ & None \\
$11 \sim 13$ & $10^4$ & $2 \times 10^4$ & $100$ & $a_i = 0$ \\
$14 \sim 15$ & $10^4$ & $2 \times 10^4$ & $100$ & $u_i < v_i$ \\
$16 \sim 20$ & $10^4$ & $2 \times 10^4$ & $100$ & None \\
\bottomrule
\end{tabular}
\end{center}

\section*{Source}

Problems of CSP-J 2023 (CCF CSP-J/S 2023 second round, junior level,
2023-10-21) as hosted on QOJ, in contest order:
\begin{itemize}
\item Problem 1: \texttt{https://qoj.ac/problem/7809}
\item Problem 2: \texttt{https://qoj.ac/problem/7810}
\item Problem 3: \texttt{https://qoj.ac/problem/7811}
\item Problem 4: \texttt{https://qoj.ac/problem/7812}
\end{itemize}
The translation follows the official Chinese statement PDF embedded in each
QOJ problem page (``Statement'' tab);
the QOJ English statement pages (\texttt{/statement/en}) were used as a
cross-check.
In the original tables, equal consecutive cells are merged;
here every row is filled in.

\end{document}
